\backmatter

\bmhead{List of abbreviations}
AI: artificial intelligence; BalAcc: balanced accuracy; CI: confidence interval; CNN--BiLSTM: convolutional neural network with bidirectional long short-term memory; CPU: central processing unit; ECE: expected calibration error; EEG: electroencephalography; EOG: electrooculography; ET: eye tracking; GAT: graph attention network; GPU: graphics processing unit; GRU: gated recurrent unit; IG: integrated gradients; INS-HDGS-CMT: Interpretable Neuro-Symbolic Hybrid Dynamic Graph Spiking Cross-Modal Transformer; LIF: leaky integrate-and-fire; LOSOCV: leave-one-subject-out cross-validation; LSTM: long short-term memory; Macro-F1: macro-averaged F1-score; MCC: Matthews correlation coefficient; MLP: multilayer perceptron; MMD: maximum mean discrepancy; PLV: phase-locking value; PR-AUC: area under the precision--recall curve; ROC-AUC: area under the receiver operating characteristic curve; ROI: region of interest; SD: standard deviation.

\bmhead{Supplementary information}
The online version contains supplementary material.

\noindent\textbf{Additional file~1:} \textbf{Note~S1.} Full specification of the implemented architecture and training objective, with every equation. \textbf{Figure~S1.} Subject-wise LOSOCV performance. \textbf{Figure~S2.} Regional EEG signatures for the case-study exemplars (S24 HIGH, S30 LOW): band power by region (frontal vs.\ posterior), the off-diagonal functional-connectivity matrix, and the within-subject HIGH$-$LOW band-power difference for frontal theta and posterior alpha. \textbf{Figure~S3.} ROI saliency vector ($5\times2$ grid) of the two case-study exemplars. \textbf{Table~S1.} EEG branch under both labelling regimes (robustness check). \textbf{Table~S2.} Top integrated-gradients input attributions. \textbf{Table~S3.} Connectivity-measure comparison ($N=16$, reduced scale). \textbf{Table~S4.} Summary of the mathematical notation used in the Methods of the main text, in order of first use. \textbf{Table~S5.} Composition of the engagement label (full weighted EEG+gaze composite definition). \textbf{Table~S6.} Gaze-free variant versus each baseline on balanced accuracy, MCC and ROC-AUC (Holm-corrected Wilcoxon, 37 folds). \textbf{Figure~S4.} Connectivity-threshold sensitivity (graph density and low-degree electrodes versus $\tau$). \textbf{Figure~S5.} Subject-count learning curve. \textbf{Figure~S6.} ROI-grid resolution sensitivity. \textbf{Figure~S7.} Activation-attribution maps of the eight rules of the rule-only variant over electrodes and frequency bands. \textbf{Figure~S8.} Per-subject scatter of the decoder against the gaze-free variant and the tuned ET-LSTM. \textbf{Table~S7.} Connectivity-threshold sensitivity (structure and downstream metrics). \textbf{Table~S8.} Baselines under a common fixed budget. \textbf{Table~S9.} Measured cost of the spiking front-end on CPU and GPU. \textbf{Table~S10.} Decoder versus each eye-tracking and fusion baseline (Holm-corrected Wilcoxon). \textbf{Table~S11.} ROI-grid sensitivity (saliency vector and model). \textbf{Table~S12.} Label-free per-subject thresholds, transductive and causal. \textbf{Table~S13.} Fidelity of the soft-rule pathway. \textbf{Table~S14.} Contribution of the eye-tracking pathway. \textbf{Table~S15.} Linear recoverability of the label from its own defining terms. \textbf{Table~S16.} Positive control (resting state versus browsing) and product-selection label, with EEGNet, ShallowConvNet, ET-LSTM and the two late fusions on the product-selection label. \textbf{Table~S17.} Fold-wise-label robustness of the decoder, its gaze-free variant, EEGNet, ShallowConvNet and ET-LSTM.

\section*{Declarations}

\bmhead{Availability of data and materials}
The dataset analysed during the current study is the publicly available NeuMa dataset described by Georgiadis et al.\ \cite{ref35}, which is deposited in the figshare repository under a CC BY 4.0 licence as a raw release, \url{https://doi.org/10.6084/m9.figshare.22117001} \cite{ref59}, and a preprocessed release, \url{https://doi.org/10.6084/m9.figshare.22117124} \cite{ref60}. No new participant recordings were collected; the derived epochs, constructed labels and analysis outputs were generated from NeuMa and are released with the code. The code supporting the findings of this study is available at \url{https://github.com/priyadharshini-D29/INS-HDGS-CMT}. All experiments use a base random seed (seed~$=42$, Table~\ref{tab2}) applied to the Python, NumPy, and PyTorch generators, with each fold's prediction averaged over a 5-member ensemble; hyperparameters, optimiser configuration, and training schedule are listed in full in Table~\ref{tab2}; software versions (Python 3.12, PyTorch 2.7.1, CUDA 12.6) are as stated in Section~\ref{sec:system}.

\bmhead{Competing interests}
The authors declare that they have no competing interests.

\bmhead{Funding}
No specific funding was received for this work. GPU computing time was provided in kind through the NVIDIA Academic Hardware Grant Program (see Acknowledgements).

\bmhead{Authors' contributions}
PD conceived and designed the study, developed the methodology and software, curated the data, performed the formal analysis and experiments, prepared the figures, and wrote the original draft. SS contributed to the conceptualisation, supervised the work, administered the project, acquired the resources and computing support, and reviewed and edited the manuscript. All authors read and approved the final manuscript.

\bmhead{Acknowledgements}
The authors thank Vellore Institute of Technology, Chennai, for research support. GPU compute for all primary experiments was provided through the NVIDIA Academic Hardware Grant via Brev.dev (NVIDIA A100-SXM4-80 GB), awarded to SS.

\bmhead{Authors' information}
Not applicable.

\bmhead{Ethics approval and consent to participate}
The NeuMa data were collected in accordance with the Declaration of Helsinki under the ethical approval reported in the original dataset publication \cite{ref35}; the present study constitutes a secondary analysis.

\bmhead{Consent for publication}
Not applicable.

\bmhead{Use of generative artificial intelligence}
Generative artificial intelligence (AI) was used solely for language editing and grammar improvement during manuscript preparation. All content was reviewed and approved by the authors, who are fully responsible for the final manuscript.

\begin{thebibliography}{60}
\small
\setlength{\itemsep}{-1.5pt}\setlength{\parsep}{0pt}

\bibitem{ref1} Alsharif AH, Mohd Isa S (2025) Electroencephalography studies on marketing stimuli: a literature review and future research agenda. International Journal of Consumer Studies 49(1):70015. \url{https://doi.org/10.1111/ijcs.70015}

\bibitem{ref2} H\"ofling TTA, Alpers GW (2023) Automatic facial coding predicts self-report of emotion, advertisement and brand effects elicited by video commercials. Frontiers in Neuroscience 17:1125983. \url{https://doi.org/10.3389/fnins.2023.1125983}

\bibitem{ref3} Gheorghe CM, Purc\u{a}rea VL, Gheorghe IR (2023) Using eye-tracking technology in neuromarketing. Romanian Journal of Ophthalmology 67(1):2-6. \url{https://doi.org/10.22336/rjo.2023.2}

\bibitem{ref6} Rawnaque FS, Rahman KM, Anwar SF, Vaidyanathan R, Chau T, Sarker F, Mamun KAA (2020) Technological advancements and opportunities in neuromarketing: a systematic review. Brain Informatics 7(1):10. \url{https://doi.org/10.1186/s40708-020-00109-x}

\bibitem{ref7} Khondakar MFK, Sarowar MH, Chowdhury MH, Majumder S, Hossain MA, Dewan MAA, Hossain QD (2024) A systematic review on EEG-based neuromarketing: recent trends and analyzing techniques. Brain Informatics 11(1):17. \url{https://doi.org/10.1186/s40708-024-00229-8}

\bibitem{ref8} Zhu L, Lv J (2023) Review of studies on user research based on EEG and eye tracking. Applied Sciences 13(11):6502. \url{https://doi.org/10.3390/app13116502}

\bibitem{ref11} Saxena S, Fink LK, Lange EB (2024) Deep learning models for webcam eye tracking in online experiments. Behavior Research Methods 56(4):3487-3503. \url{https://doi.org/10.3758/s13428-023-02190-6}

\bibitem{ref12} Panda D, Das Chakladar D, Rana S, Shamsudin MN (2024) Spatial attention-enhanced EEG analysis for profiling consumer choices. IEEE Access 12:13477-13487. \url{https://doi.org/10.1109/ACCESS.2024.3355977}

\bibitem{ref13} Khurana V, Gahalawat M, Kumar P, Roy PP, Dogra DP, Scheme E, Soleymani M (2021) A survey on neuromarketing using EEG signals. IEEE Transactions on Cognitive and Developmental Systems 13(4):732-749. \url{https://doi.org/10.1109/TCDS.2021.3065200}

\bibitem{ref18} Kalaganis FP, Georgiadis K, Oikonomou VP, Laskaris NA, Nikolopoulos S, Kompatsiaris I (2025) A hybrid neuromarketing approach exploiting EEG graph signal processing and gaze dynamic patterning. Brain Informatics 12(1):23. \url{https://doi.org/10.1186/s40708-025-00272-z}

\bibitem{ref19} Lin X, Chen J, Ma W, Tang W, Wang Y (2023) EEG emotion recognition using improved graph neural network with channel selection. Computer Methods and Programs in Biomedicine 231:107380. \url{https://doi.org/10.1016/j.cmpb.2023.107380}

\bibitem{ref20} Chen W, Liao Y, Dai R, Dong Y, Huang L (2024) EEG-based emotion recognition using graph convolutional neural network with dual attention mechanism. Frontiers in Computational Neuroscience 18:1416494. \url{https://doi.org/10.3389/fncom.2024.1416494}

\bibitem{ref21} Yan H, Guo K, Xing X, Xu X (2024) Bridge graph attention based graph convolution network with multi-scale transformer for EEG emotion recognition. IEEE Transactions on Affective Computing 15(4):2042-2054. \url{https://doi.org/10.1109/TAFFC.2024.3394873}

\bibitem{ref22} Jin M, Du C, He H, Cai T, Li J (2024) PGCN: pyramidal graph convolutional network for EEG emotion recognition. IEEE Transactions on Multimedia 26:9070-9082. \url{https://doi.org/10.1109/TMM.2024.3385676}

\bibitem{ref24} Huang SC, Pareek A, Seyyedi S, Banerjee I, Lungren MP (2020) Fusion of medical imaging and electronic health records using deep learning: a systematic review and implementation guidelines. npj Digital Medicine 3(1):136. \url{https://doi.org/10.1038/s41746-020-00341-z}

\bibitem{ref25} Fu B, Gu C, Fu M, Xia Y, Liu Y (2023) A novel feature fusion network for multimodal emotion recognition from EEG and eye movement signals. Frontiers in Neuroscience 17:1234162. \url{https://doi.org/10.3389/fnins.2023.1234162}

\bibitem{ref26} Jin X, Xiao J, Jin L, Zhang X (2024) Residual multimodal transformer for expression--EEG fusion continuous emotion recognition. CAAI Transactions on Intelligence Technology 9(5):1290-1304. \url{https://doi.org/10.1049/cit2.12346}

\bibitem{ref27} Pillalamarri R, Shanmugam U (2025) A review on EEG-based multimodal learning for emotion recognition. Artificial Intelligence Review 58(5):131. \url{https://doi.org/10.1007/s10462-025-11126-9}

\bibitem{ref28} Gu Y, Zhong X, Qu C, Liu C, Chen B (2023) A domain generative graph network for EEG-based emotion recognition. IEEE Journal of Biomedical and Health Informatics 27(5):2377-2386. \url{https://doi.org/10.1109/JBHI.2023.3242090}

\bibitem{ref29} Wu X, Ju X, Dai S, Li X, Li M (2024) Multi-source domain adaptation for EEG emotion recognition based on inter-domain sample hybridization. Frontiers in Human Neuroscience 18:1464431. \url{https://doi.org/10.3389/fnhum.2024.1464431}

\bibitem{ref30} Lu W, Liu H, Ma H, Tan TP, Xia L (2023) Hybrid transfer learning strategy for cross-subject EEG emotion recognition. Frontiers in Human Neuroscience 17:1280241. \url{https://doi.org/10.3389/fnhum.2023.1280241}

\bibitem{ref31} Liang S, Li L, Zu W, Feng W, Hang W (2024) Adaptive deep feature representation learning for cross-subject EEG decoding. BMC Bioinformatics 25(1):393. \url{https://doi.org/10.1186/s12859-024-06024-w}

\bibitem{ref32} Xu C, Song Y, Zheng Q, Wang Q, Heng PA (2025) Unsupervised multi-source domain adaptation via contrastive learning for EEG classification. Expert Systems with Applications 261:125452. \url{https://doi.org/10.1016/j.eswa.2024.125452}

\bibitem{ref33} Hakim A, Klorfeld S, Sela T, Friedman D, Shabat-Simon M, Levy DJ (2021) Machines learn neuromarketing: improving preference prediction from self-reports using multiple EEG measures and machine learning. International Journal of Research in Marketing 38(3):770-791. \url{https://doi.org/10.1016/j.ijresmar.2020.10.005}

\bibitem{ref35} Georgiadis K, Kalaganis FP, Riskos K, Matta E, Oikonomou VP, Yfantidou I, Chantziaras D, Pantouvakis K, Nikolopoulos S, Laskaris NA, et al (2023) NeuMa -- the absolute neuromarketing dataset en route to an holistic understanding of consumer behaviour. Scientific Data 10(1):508. \url{https://doi.org/10.1038/s41597-023-02392-9}

\bibitem{ref36} Usman SM, Khalid S, Tanveer A, Imran AS, Zubair M (2025) Multimodal consumer choice prediction using EEG signals and eye tracking. Frontiers in Computational Neuroscience 18:1516440. \url{https://doi.org/10.3389/fncom.2024.1516440}

\bibitem{ref38} Burkitt AN (2006) A review of the integrate-and-fire neuron model: I. homogeneous synaptic input. Biological Cybernetics 95(1):1-19. \url{https://doi.org/10.1007/s00422-006-0068-6}

\bibitem{ref39} Zenke F, Ganguli S (2018) Superspike: supervised learning in multilayer spiking neural networks. Neural Computation 30(6):1514-1541. \url{https://doi.org/10.1162/neco_a_01086}

\bibitem{ref40} Cai S, Lin Z, Liu X, Wei W, Wang S, Zhang M, Schultz T, Li H (2026) Spiking neural networks for EEG signal analysis: from theory to practice. Neural Networks 194:108127. \url{https://doi.org/10.1016/j.neunet.2025.108127}

\bibitem{ref41} Li W, Fang C, Zhu Z, Chen C, Song A (2024) Fractal spiking neural network scheme for EEG-based emotion recognition. IEEE Journal of Translational Engineering in Health and Medicine 12:106-118. \url{https://doi.org/10.1109/JTEHM.2023.3320132}

\bibitem{ref42} Xu F, Pan D, Zheng H, Ouyang Y, Jia Z, Zeng H (2024) EESCN: a novel spiking neural network method for EEG-based emotion recognition. Computer Methods and Programs in Biomedicine 243:107927. \url{https://doi.org/10.1016/j.cmpb.2023.107927}

\bibitem{ref43} Veli\v{c}kovi\'c P, Cucurull G, Casanova A, Romero A, Li\`o P, Bengio Y (2018) Graph attention networks. In: International Conference on Learning Representations (ICLR). \url{https://openreview.net/forum?id=rJXMpikCZ}

\bibitem{ref44} Vaswani A, Shazeer N, Parmar N, Uszkoreit J, Jones L, Gomez AN, Kaiser L, Polosukhin I (2017) Attention is all you need. In: Advances in Neural Information Processing Systems, vol 30. \url{https://proceedings.neurips.cc/paper/2017/hash/3f5ee243547dee91fbd053c1c4a845aa-Abstract.html}

\bibitem{ref45} Rudin C (2019) Stop explaining black box machine learning models for high stakes decisions and use interpretable models instead. Nature Machine Intelligence 1(5):206-215. \url{https://doi.org/10.1038/s42256-019-0048-x}

\bibitem{ref46} Koh PW, Nguyen T, Tang YS, Mussmann S, Pierson E, Kim B, Liang P (2020) Concept bottleneck models. In: International Conference on Machine Learning (ICML), vol 119, pp 5338-5348. \url{https://proceedings.mlr.press/v119/koh20a.html}

\bibitem{ref47} Garcez Ad, Lamb LC (2023) Neurosymbolic AI: the 3rd wave. Artificial Intelligence Review 56(11):12387-12406. \url{https://doi.org/10.1007/s10462-023-10448-w}

\bibitem{ref48} Lin TY, Goyal P, Girshick R, He K, Doll\'ar P (2017) Focal loss for dense object detection. In: Proceedings of the IEEE International Conference on Computer Vision (ICCV), pp 2980-2988. \url{https://doi.org/10.1109/ICCV.2017.324}

\bibitem{ref49} Cui Y, Jia M, Lin TY, Song Y, Belongie S (2019) Class-balanced loss based on effective number of samples. In: Proceedings of the IEEE/CVF Conference on Computer Vision and Pattern Recognition (CVPR), pp 9268-9277. \url{https://doi.org/10.1109/CVPR.2019.00949}

\bibitem{ref50} Lawhern VJ, Solon AJ, Waytowich NR, Gordon SM, Hung CP, Lance BJ (2018) EEGNet: a compact convolutional neural network for EEG-based brain--computer interfaces. Journal of Neural Engineering 15(5):056013. \url{https://doi.org/10.1088/1741-2552/aace8c}

\bibitem{ref51} Schirrmeister RT, Springenberg JT, Fiederer LDJ, Glasstetter M, Eggensperger K, Tangermann M, Hutter F, Burgard W, Ball T (2017) Deep learning with convolutional neural networks for EEG decoding and visualization. Human Brain Mapping 38(11):5391-5420. \url{https://doi.org/10.1002/hbm.23730}

\bibitem{ref52} Ding Y, Robinson N, Zhang S, Zeng Q, Guan C (2023) TSception: capturing temporal dynamics and spatial asymmetry from EEG for emotion recognition. IEEE Transactions on Affective Computing 14(3):2238-2250. \url{https://doi.org/10.1109/TAFFC.2022.3169001}

\bibitem{ref54} Sundararajan M, Taly A, Yan Q (2017) Axiomatic attribution for deep networks. In: Proceedings of the 34th International Conference on Machine Learning, vol 70, pp 3319-3328. \url{https://proceedings.mlr.press/v70/sundararajan17a.html}

\bibitem{ref55} Jiang WB, Zhao LM, Lu BL (2024) Large brain model for learning generic representations with tremendous EEG data in BCI. In: International Conference on Learning Representations (ICLR). \url{https://openreview.net/forum?id=QzTpTRVtrP}

\bibitem{ref56} Wang G, Liu W, He Y, Xu C, Ma L, Li H (2024) EEGPT: pretrained transformer for universal and reliable representation of EEG signals. In: Advances in Neural Information Processing Systems, vol 37. \url{https://openreview.net/forum?id=lvS2b8CjG5}

\bibitem{ref57} Aziz MZ, Yu X, Guo X, He X, Huang B, Fan Z (2025) BCINetV1: integrating temporal and spectral focus through a novel convolutional attention architecture for MI EEG decoding. Sensors 25(15):4657. \url{https://doi.org/10.3390/s25154657}

\bibitem{ref58} Huang B, Aziz MZ, Guo X, He X, Zheng J, Yu X (2026) A hybrid adaptive EEG signal tokenization and vector embedding framework for interpretable brain--computer interfaces. Biomedical Signal Processing and Control 126:110765. \url{https://doi.org/10.1016/j.bspc.2026.110765}

\bibitem{ref59} Georgiadis K, Kalaganis FP, Riskos K, Matta E, Oikonomou VP, Yfantidou I, Chantziaras D, Pantouvakis K, Nikolopoulos S, Laskaris NA, et al (2023) NeuMa (Raw): a multimodal neuromarketing dataset. figshare. Dataset. \url{https://doi.org/10.6084/m9.figshare.22117001}

\bibitem{ref60} Georgiadis K, Kalaganis FP, Riskos K, Matta E, Oikonomou VP, Yfantidou I, Chantziaras D, Pantouvakis K, Nikolopoulos S, Laskaris NA, et al (2023) NeuMa (PreProcessed): a multimodal neuromarketing dataset. figshare. Dataset. \url{https://doi.org/10.6084/m9.figshare.22117124}

